\scanpage{45}
autre que $\underline b(\Omega)$ du \S~5.

Fixons un $\underline b(q)$ \inserted{$=\mathrm{tr}$} i.e. \deleted{\unsure{$q$}} \uline{la forme} $\mathrm{tr}$
\uline{trace} $B\longrightarrow\mathcal O_S$, et étudions les $q$ correspondants
à la forme trace donnée. On trouve
que ceux-ci \deleted{\unsure{font un|faisceau}} \inserted{sont les sections d'un faisceau qui est un} \deleted{\unsure{pré}}torseur (et $=$ un
torseur\ldots) \unsure{sous|sur|à} le faisceau ou Module des
formes quadratiques $q_0$ telles que
\[
\left\{\begin{aligned}q_0(e_B)&=0\\
\varphi_0(e_B,\unsure{-|\cdot})&=0\qquad\deleted{\forall x}
\end{aligned}\right.\quad
\substack{\text{(\unsure{où|ici} $\varphi_0$ est la}\\\text{forme bil associée}\\\text{à $q$ )}} \tag*{(82)}
\]
\deleted{l'hom $V\longrightarrow\check V$}
\deleted{i.e. $\varphi_0(e_B,\unsure{-})$ une forme $\varphi_0$}
ce sont aussi les formes quadratiques
qui se factorisent par $B\longrightarrow\check L$, i.e.
qui correspondent aux formes quadratiques
sur $L$, ou encore aux \deleted{\unsure{fo|for}} sections de
$L^{\otimes2}$. Donc les formes en question
forment \deleted{les} sections\deleted{\unsure{en}} d'un fibré affine sous
$L^{\otimes2}$, dont les sections $c_0$ correspondent aux $q_0$ par
\[
q_0(x)=c_0\,\pi_1(x)^2. \tag*{(83)}
\]
% archive sheet number: 44
